\chapter{Local Data, Spectral Measures, and Green Functions}
\label{ch:volume-ii-local-data}

\section{Local Data for Scattering and Renormalization}

The constructions in this part use the local Hilbert-space data already
defined. Fix a spacetime dimension \(D\), write \(d=D-1\), and use the
mostly-plus metric
\[
  \eta_{\mu\nu}=\operatorname{diag}(-,+,\ldots,+),
  \qquad
  p^2=\eta_{\mu\nu}p^\mu p^\nu .
\]
The quantum data are:
\begin{enumerate}
  \item a Hilbert space \(\Hilb\);
  \item a strongly continuous unitary representation
  \(U(a,\Lambda)\) of the proper orthochronous Poincare group;
  \item commuting self-adjoint translation generators \(P^\mu\), defined by
  \[
    U(a,\mathbf 1)=\ee^{\ii a_\mu P^\mu};
  \]
  \item the spectrum condition
  \[
    \operatorname{Spec}(P)\subset\overline V_+
    =
    \{p\in\mathbb R^D:p^0\ge0,\ p^2\le0\};
  \]
  \item a Poincare-invariant unit vector \(\Omega\in\Hilb\);
  \item for each bounded spacetime region \(O\), a local observable algebra
  \(\Obs(O)\), or equivalently a compatible family of smeared fields on a
  common invariant domain.
\end{enumerate}
Locality means that if \(O_1\) and \(O_2\) are spacelike separated, then
bounded observables \(A_i\in\Obs(O_i)\) commute. For fields, the corresponding
statement is imposed after smearing with test functions whose supports are
spacelike separated.

These data support several constructions, each with its own hypotheses:
\[
\begin{tikzpicture}[>=Stealth,scale=0.92,
  box/.style={draw,thick,align=center,minimum width=3.05cm,
    minimum height=0.92cm,fill=blue!4},
  arr/.style={->,thick}]
  \node[box] (local) at (0,0) {local Hilbert-space\\data};
  \node[box] (spectral) at (4.15,0) {spectral\\measures};
  \node[box] (greens) at (8.3,0) {Green-function\\distributions};
  \node[box,fill=green!5] (scatter) at (4.15,-1.85)
    {asymptotic states\\and kernels};
  \node[box,fill=orange!7] (ren) at (8.3,-1.85)
    {renormalized\\continuum limits};
  \draw[arr] (local)--(spectral);
  \draw[arr] (spectral)--(greens);
  \draw[arr] (spectral)--(scatter);
  \draw[arr] (greens)--(scatter);
  \draw[arr] (greens)--(ren);
\end{tikzpicture}
\]
The first arrow is the joint spectral theorem applied to the translation
generators. The arrows to scattering use isolated one-particle spectral
support and large-time limits. The arrow to renormalized continuum limits uses
regulated correlation functions and the locality of counterterms.

\section{Spectral Measures}

Let \(E(\Delta)\) be the joint spectral projection of \(P^\mu\) for a Borel
set \(\Delta\subset\mathbb R^D\). If \(\widehat\Phi_\alpha(f)\) is a smeared
local field and \(v_\alpha=\widehat\Phi_\alpha(f)\Omega\), the positive
measure
\[
  \nu_{\Phi,f}(\Delta)
  =
  \langle v_\alpha,E(\Delta)v_\alpha\rangle_{\Hilb}
\]
records the energy-momentum content created from the vacuum by that smeared
field. For a scalar field \(\widehat\phi\), translation invariance and Lorentz
covariance allow the two-point Wightman distribution to be written in the
K\"all\'en--Lehmann form
\[
  W_2(x-y)
  =
  \langle\Omega|\widehat\phi(x)\widehat\phi(y)|\Omega\rangle
  =
  \int_0^\infty \dd\rho_\phi(\mu^2)\,
  \Delta_+(x-y;\mu^2),
\]
where \(\dd\rho_\phi\) is a positive Borel measure and
\[
  \Delta_+(x;\mu^2)
  =
  \int\frac{\dd^D k}{(2\pi)^{D-1}}\,
  \theta(k^0)\delta(k^2+\mu^2)\,\ee^{\ii k\cdot x}.
\]
The measure \(\dd\rho_\phi\) depends on the chosen interpolating field, while
the spectral support of the full theory is an invariant property of the
translation representation.

If a stable scalar particle of mass \(m\) is created by \(\widehat\phi\), the
measure contains an atom
\[
  \dd\rho_\phi(\mu^2)
  =
  Z_\phi\,\delta(\mu^2-m^2)\,\dd\mu^2
  +
  \dd\rho_{\phi,\mathrm c}(\mu^2),
  \qquad
  Z_\phi>0.
\]
The continuous part \(\dd\rho_{\phi,\mathrm c}\) begins at the lowest invariant
mass threshold accessible to \(\widehat\phi\).

\begin{center}
\begin{tikzpicture}[>=Stealth,scale=0.98]
  \draw[->,line width=0.7pt] (0,0)--(7.0,0) node[right] {$\mu^2$};
  \draw[->,line width=0.7pt] (0,0)--(0,3.1) node[above] {$\dd\rho_\phi$};
  \draw[blue!75!black,line width=1.3pt] (1.85,0)--(1.85,2.78);
  \node[blue!75!black,below] at (1.85,-0.1) {$m^2$};
  \node[blue!75!black,anchor=west] at (2.08,2.78)
    {$Z_\phi\delta(\mu^2-m^2)$};
  \draw[red!75!black,line width=1.4pt]
    (3.65,0) .. controls (3.9,0.92) and (4.45,1.45) .. (5.2,1.28)
    .. controls (5.85,1.13) and (6.25,0.73) .. (6.65,0.58);
  \node[red!75!black,below] at (3.65,-0.1) {$M_{\mathrm{thr}}^2$};
  \node[red!75!black,align=center] at (5.75,2.25)
    {continuum\\support};
  \draw[red!75!black,->,line width=0.7pt] (5.55,1.82)--(5.08,1.23);
  \draw[decorate,decoration={brace,amplitude=5pt},gray!75]
    (1.98,0.56)--(3.48,0.56);
  \node[gray!70!black] at (2.75,0.96) {gap};
\end{tikzpicture}
\end{center}

This picture is the input for later scattering analysis. The atom supplies a
one-particle Hilbert space and pole residue. The continuum supplies threshold
data, branch cuts of Green functions, and the channels in which bound states
or resonances may appear.

\section{Lorentzian and Euclidean Green Functions}

For local fields \(\widehat\Phi_{a_1},\ldots,\widehat\Phi_{a_N}\), where
\(a_i\) denotes all tensor, spinor, and species labels, define the
time-ordered distribution
\[
  G_{a_1\cdots a_N}(x_1,\ldots,x_N)
  =
  \langle\Omega|
  \operatorname T\{
    \widehat\Phi_{a_1}(x_1)\cdots
    \widehat\Phi_{a_N}(x_N)
  \}
  |\Omega\rangle .
\]
The expression is a distributional matrix element. Products at coincident
points require a specified time-ordering prescription or a renormalized
definition of composite insertions.

For the scalar two-point function, the spectral representation gives
\[
  G_2(x)
  =
  \int_0^\infty \dd\rho_\phi(\mu^2)\,\Delta_F(x;\mu^2),
\]
with
\[
  \Delta_F(x;\mu^2)
  =
  -\ii
  \int\frac{\dd^D k}{(2\pi)^D}\,
  \frac{\ee^{\ii k\cdot x}}{k^2+\mu^2-\ii0}.
\]
The \(+\ii0\) or \(-\ii0\) notation is a boundary-value prescription. With
the mostly-plus convention above, the displayed denominator is the Feynman
boundary value of the analytic function away from the mass shell.

Euclidean points are written \(x_E=(\vec x,\tau)\in\mathbb R^D\), with
\(x^0=-\ii\tau\) in the standard forward rotation. In a framework where
Schwinger functions \(S_N\) have been constructed, the scalar two-point
Schwinger function has the representation
\[
  S_2(x_E)
  =
  \int_0^\infty \dd\rho_\phi(\mu^2)\,\Delta_E(x_E;\mu^2),
\]
where
\[
  \Delta_E(x_E;\mu^2)
  =
  \int\frac{\dd^D k_E}{(2\pi)^D}\,
  \frac{\ee^{\ii k_E\cdot x_E}}{k_E^2+\mu^2}.
\]
The passage between Wightman, time-ordered, and Euclidean functions is
controlled by spectral support, analyticity domains, and reconstruction
conditions. In perturbation theory it is implemented as a contour and
boundary-value calculation; in an Euclidean constructive framework it is part
of the reconstruction problem.

\section{Regulated Path-Integral Calculus}

Let \(\Lambda\) be a regulator scale. A regulated Euclidean field
configuration space is denoted by \(\mathcal C_\Lambda\), and its integration
symbol by \(D_\Lambda\varphi\). The regulated action is a functional
\[
  S_{\Lambda,E}[\varphi;g_I(\Lambda)]
  =
  \sum_I g_I(\Lambda) O_{I,\Lambda}[\varphi],
\]
where \(O_{I,\Lambda}\) are regulated local or quasi-local functionals and
\(g_I(\Lambda)\) are cutoff-dependent coefficients. For a source \(J\), the
regulated generating functional is
\[
  Z_\Lambda[J]
  =
  \int_{\mathcal C_\Lambda}D_\Lambda\varphi\,
  \exp\left(
    -S_{\Lambda,E}[\varphi;g_I(\Lambda)]
    +\int \dd^D x_E\,J(x_E)\varphi(x_E)
  \right).
\]
Regulated \(N\)-point functions are obtained by differentiating
\[
  \langle\varphi(x_1)\cdots\varphi(x_N)\rangle_\Lambda
  =
  \frac{1}{Z_\Lambda[0]}
  \frac{\delta^N Z_\Lambda[J]}
  {\delta J(x_1)\cdots\delta J(x_N)}
  \bigg|_{J=0}.
\]

A continuum Schwinger function is a limit
\[
  S_N(x_1,\ldots,x_N)
  =
  \lim_{\Lambda\to\infty}
  \langle\varphi(x_1)\cdots\varphi(x_N)\rangle_\Lambda
\]
in a specified distribution topology, after choosing the cutoff dependence of
the coefficients \(g_I(\Lambda)\). Existence of this limit is a dynamical
claim. Perturbative renormalization proves the existence of finite formal
coefficients order by order in selected classes of theories; constructive
methods prove stronger nonperturbative statements in selected models.

\begin{center}
\begin{tikzpicture}[>=Stealth,scale=0.94,
  box/.style={draw,thick,align=center,minimum width=3.25cm,
    minimum height=0.86cm,fill=green!5},
  arr/.style={->,thick}]
  \node[box] (cutoff) at (0,0)
    {cutoff data\\\(\mathcal C_\Lambda,S_{\Lambda,E}\)};
  \node[box] (coeff) at (4.35,0)
    {local coefficients\\\(g_I(\Lambda)\)};
  \node[box] (corr) at (8.7,0)
    {regulated\\correlators};
  \node[box,fill=orange!7] (limit) at (4.35,-1.85)
    {distributional\\continuum limit};
  \draw[arr] (cutoff)--(coeff);
  \draw[arr] (coeff)--(corr);
  \draw[arr] (corr)--(limit);
  \draw[arr] (coeff)--(limit);
\end{tikzpicture}
\end{center}

The Lorentzian functional integral
\[
  \int D\phi\,
  \exp\left(\ii S[\phi]+\ii\int\dd^D x\,J(x)\phi(x)\right)
\]
is used in this monograph as an oscillatory or perturbative representation of
time-ordered Green functions after a regulator and a boundary prescription
have been specified.

\section{Three Distinct Operations}

The later chapters repeatedly use three operations whose domains are distinct.

First, a Green-function expansion is an expansion of the regulated or formal
generating functional. Its output is a collection of distributions
\[
  G_N=G_N^{(0)}+G_N^{(1)}+\cdots
\]
with external momenta left off shell until a boundary value or residue is
taken.

Second, LSZ reduction is a pole-residue operation. In a stable scalar sector
with one-particle residue \(Z_\phi\), it applies factors
\[
  Z_\phi^{-1/2}\,\ii(k^2+m^2)
\]
to external legs of a connected time-ordered Green function and then takes
the distributional boundary value on the positive-energy mass shells.

Third, renormalization is the organization of cutoff dependence in
\(g_I(\Lambda)\), field normalizations, and composite-operator definitions so
that selected correlation functions or matrix elements have finite limits.
The locality of the counterterms is the structural fact that connects
renormalization to the original local data.

The sequence in this part is therefore fixed: spectral data identify stable
one-particle inputs and thresholds; scattering kernels are defined from
asymptotic states; LSZ connects those kernels to Green functions; analytic and
renormalization structures are then studied as properties of these already
defined objects.
